% =========================================================================
% Figura del problema 2.2: (a) carga q en (a,b,c); (b) tras añadir -q en el
% origen, el sistema es un dipolo de momento p = q r'_1. Perspectiva tdplot.
% =========================================================================
\begin{figure}[ht]
\centering
\newcommand{\ejesIII}{%
  \draw[eje] (0,0,0) -- (2.6,0,0) node[anchor=north east, etiq] {$x$};
  \draw[eje] (0,0,0) -- (0,2.7,0) node[anchor=west, etiq] {$y$};
  \draw[eje] (0,0,0) -- (0,0,2.3) node[anchor=south, etiq] {$z$};}
\newcommand{\proyecciones}{%
  \draw[punteado] (1.6,1.9,0) -- (1.6,1.9,1.5);
  \draw[punteado] (1.6,0,0) -- (1.6,1.9,0) -- (0,1.9,0);
  \node[etiq, font=\footnotesize, anchor=north] at (1.6,0,0) {$a$};
  \node[etiq, font=\footnotesize, anchor=north west] at (0,1.9,0) {$b$};
  \draw[punteado] (0,0,1.5) -- (1.6,1.9,1.5);
  \node[etiq, font=\footnotesize, anchor=east] at (0,0,1.5) {$c$};}
\tdplotsetmaincoords{70}{120}
\begin{minipage}[b]{0.47\textwidth}\centering
\begin{tikzpicture}[tdplot_main_coords, scale=1.25, >=Stealth]
  \ejesIII \proyecciones
  \draw[vect, rojoteo] (0,0,0) -- (1.6,1.9,1.5) node[pos=0.55, anchor=south east, etiq, text=rojoteo] {$\rp_1$};
  \node[cargapos] at (1.6,1.9,1.5) {$+$};
  \node[etiq, anchor=south west] at (1.6,1.9,1.62) {$q\ (a,b,c)$};
\end{tikzpicture}\\[2pt]{\small\scshape (a)}
\end{minipage}\hfill
\begin{minipage}[b]{0.47\textwidth}\centering
\begin{tikzpicture}[tdplot_main_coords, scale=1.25, >=Stealth]
  \ejesIII \proyecciones
  \draw[->, line width=1.5pt, magacero] (0,0,0) -- (1.6,1.9,1.5) node[pos=0.45, anchor=south east, xshift=-3pt, yshift=4pt, etiq, text=magacero] {$\vp=q\,\rp_1$};
  \node[cargapos] at (1.6,1.9,1.5) {$+$};
  \node[carganeg] at (0,0,0) {$-$};
  \node[etiq, anchor=south west] at (1.6,1.9,1.62) {$q$};
  \node[etiq, anchor=north west] at (0,0.1,-0.15) {$-q$};
\end{tikzpicture}\\[2pt]{\small\scshape (b)}
\end{minipage}
\caption{Problema~2.2. (a)~Carga puntual $q$ en $\rp_1=(a,b,c)$: $Q=q$, $\vp=q\,\rp_1$ y $\Qjk=q\,(3x'_{1,j}x'_{1,k}-\delt{j}{k}r_1'^{\,2})$. (b)~Al añadir $-q$ en el origen, el sistema se
convierte en un dipolo físico: solo cambia la carga neta ($Q=0$), mientras que $\vp$ y $\Qt$ no varían porque la nueva carga está en $\rp_2=\vc{0}$.}
\label{fig:sem-carga}
\end{figure}
